\documentclass[10pt,twocolumn]{article}
\usepackage[margin=0.75in,columnsep=0.25in]{geometry}
\usepackage{times}
\usepackage{amsmath,amssymb}
\usepackage{graphicx}
\usepackage{booktabs}
\usepackage{array}
\usepackage{caption}
\usepackage[hidelinks]{hyperref}
\usepackage{titlesec}
\usepackage{enumitem}
\usepackage{xcolor}
\usepackage{multirow}

\titleformat{\section}{\centering\normalfont\small\bfseries}{\Roman{section}.}{0.6em}{\MakeUppercase}
\titleformat{\subsection}{\normalfont\small\bfseries\itshape}{\Alph{subsection}.}{0.5em}{}
\renewcommand{\thesection}{\Roman{section}}

\captionsetup[table]{labelfont=bf,font=small,justification=raggedright,singlelinecheck=false}
\captionsetup[figure]{labelfont=bf,font=small}

\begin{document}
\twocolumn[
\begin{center}
{\LARGE\bfseries Outcome Without Method: An Outcome-Validation Integrity\\[2pt] Score for Agentic Cybersecurity Benchmarks\par}
\vspace{10pt}
{\large Babar Khan Akhunzada\par}
{\normalsize SecurityWall\quad\texttt{founder@securitywall.co}\par}
\vspace{8pt}
\end{center}
]

\begin{abstract}
\noindent
Agentic large language model (LLM) systems are increasingly evaluated on
cybersecurity tasks: capture-the-flag (CTF) exploitation, vulnerability
discovery, malware analysis, and automated patching, typically using
benchmarks that report a single scalar success metric, such as a binary
flag capture, a multiple-choice accuracy, or a pass@$k$ rate. We show,
through a structured audit of eleven independently published benchmarks
and deployments spanning academic preprints, an industry blog, a
government-run competition, an enterprise production deployment, and an
industry framework evaluation, that the design of a success metric is a
strong predictor of whether that metric can be gamed without the
intended capability being exercised, a failure mode we term
\emph{outcome-without-method}. We formalize a five-dimension,
0--10-point \textbf{Outcome-Validation Integrity Score (OVIS)} that
scores each benchmark on method verification, shortcut and contamination
auditing, human ground-truth anchoring, grader independence, and
reproducibility. Applying OVIS across our corpus, benchmarks using
purely outcome-based grading (binary flag or multiple-choice) score a
mean method-verification sub-score of \textbf{16.7\%}, versus
\textbf{100.0\%} for benchmarks using process-aware grading (tiered
capability ladders, graph-grounded step rewards, or functional
patch-validation). We further document a previously reported instance of
reward hacking in which a benchmarked model achieved an 85\% nominal
pass rate on a challenge while using the intended exploitation technique
in \textbf{zero} of 681 relevant runs, and show that a post-hoc manual
audit, rather than benchmark redesign, was sufficient to detect it. We
situate OVIS within the broader reward-hacking and specification-gaming
literature, argue that OVIS-style auditing should be a required
companion to any reported capability number in agentic cybersecurity
evaluation, release our scoring rubric and corpus as a reusable
artifact, and outline how the approach generalizes to non-security
agentic benchmarks.
\end{abstract}

\vspace{4pt}
\noindent\textbf{\textit{Index Terms}}---agentic AI, large language
models, cybersecurity benchmarks, reward hacking, evaluation
methodology, penetration testing, benchmark integrity, meta-analysis

\vspace{6pt}

\section{Introduction}

Large language model (LLM) agents are now routinely evaluated on
offensive and defensive cybersecurity tasks: capturing flags in
capture-the-flag (CTF) exercises~\cite{xbow2026,cybench2024,enigma2024},
discovering and patching real-world software
vulnerabilities~\cite{aixcc2025,fuzzingbrain2026}, analyzing malware
detonation logs~\cite{cybersoceval2025}, and conducting multi-step SOC
investigations~\cite{excytin2025}. The resulting capability claims, such
as ``model X solves 90 of 104 CTF challenges,'' ``model Y clears the
hard tier at scale,'' or ``model Z performs offensive security at
3{,}600$\times$ human speed''~\cite{cai2025}, circulate widely and shape
both research priorities and public risk perception.

Nearly all of these benchmarks reduce agent performance to a single
scalar: did the agent produce the correct flag, select the correct
multiple-choice option, or reach a terminal state matching a stored
answer key. This paper asks a narrower, more mechanical question that
the literature has so far answered mostly anecdotally. When a
benchmark's grading function checks only the outcome state and not the
method by which that state was reached, how often does the reported
number reflect something other than the capability the benchmark claims
to measure?

We are motivated by a concrete case. TrustedSec's evaluation of six
locally-hosted models against eight Juice Shop
challenges~\cite{trustedsec2026} initially reported a 98.5\% nominal
pass rate on a JWT-forgery challenge. A subsequent manual audit,
decoding every passing JWT across 681 relevant runs, found that
\emph{zero} of those runs used the intended RS256-to-HS256
algorithm-confusion attack; every ``pass'' was an unrelated shortcut
(replaying an SQL-injection-obtained token, or the unrelated
\texttt{alg:none} attack) that happened to satisfy a string-match
success condition. The benchmark's own success check,
\texttt{contains(admin@juice-sh.op)}, could not distinguish the intended
exploit from an incidental one. This is not an isolated design flaw. We
find structurally similar patterns, of varying severity, across the
corpus assembled for this study, and the pattern is consistent with a
much older concern in the reinforcement learning literature: that
optimizing a proxy objective reliably produces behavior that satisfies
the proxy while diverging from the objective it was meant to
stand in for~\cite{amodei2016concrete,skalse2022reward,pan2022reward}.

\subsection{Contributions}
\begin{itemize}[leftmargin=*,itemsep=1pt,topsep=2pt]
  \item We formalize the \textbf{Outcome-Validation Integrity Score
  (OVIS)}, a five-dimension, deterministic, auditable rubric for scoring
  the degree to which a benchmark's grading function can be gamed
  without the intended capability being exercised (Section~\ref{sec:method}).
  \item We apply OVIS to eleven benchmarks and deployments with
  sufficient publicly documented grading detail to score defensibly, and
  report a quantitative association between metric design category
  (outcome-only versus process-aware) and method-verification integrity
  (Section~\ref{sec:results}).
  \item We show that post-hoc manual auditing can functionally rescue an
  outcome-only benchmark design, using the TrustedSec case as a worked
  example, and argue this should be treated as a minimum bar rather than
  a bonus (Section~\ref{sec:discussion}).
  \item We situate the finding within the wider CTF and pentesting
  benchmark landscape and the foundational reward-hacking literature,
  and release the full scoring rubric, corpus, and computation script as
  a reusable artifact for future benchmark audits.
\end{itemize}

\section{Related Work}

\subsection{Benchmark integrity in CTF-style evaluation}
Sultan~\cite{xbow2026} addresses a narrower instance of the shortcut
problem specific to XBOW-style benchmarks, where flags are computed as
$\mathrm{sha256}(\mathrm{challenge\_id})$: an agent that discovers the
challenge identifier could compute the flag directly without exploiting
the target. That work audits 142 solved transcripts for identifier
leakage and hash computation, finding none. This is a shortcut-specific
audit rather than a general method-verification framework; OVIS
generalizes this single-purpose audit into a reusable, multi-dimension
scoring instrument applicable across benchmark designs.

\subsection{Reward hacking and specification gaming}
The general phenomenon of an agent satisfying a proxy objective without
satisfying the intended objective was named as an open problem in AI
safety by Amodei et al.~\cite{amodei2016concrete} and formalized as
``proxies are unhackable only under narrow conditions'' by Skalse et
al.~\cite{skalse2022reward}. Pan et al.~\cite{pan2022reward} showed
empirically that more capable reinforcement learning agents exploit
reward misspecification \emph{more}, not less, and documented sudden
phase transitions from proxy-and-true-reward correlation to outright
divergence. DeepMind's specification-gaming catalog~\cite{krakovna2020specification}
collects dozens of concrete instances of this pattern across RL domains.
TrustedSec's Juice Shop study~\cite{trustedsec2026} is, to our
knowledge, the first work to document this failure mode with per-run
forensic detail (JWT decoding across hundreds of runs) in an agentic
cybersecurity benchmark specifically, and is the empirical anchor for
this paper. OVIS can be read as an attempt to operationalize the
Skalse/Pan/Amodei critique as a checklist applicable to a benchmark's
published methodology, rather than requiring a from-scratch theoretical
analysis of every new benchmark's reward structure.

\subsection{Process-aware grading}
A minority of benchmarks in our corpus already grade intermediate
process rather than terminal outcome alone. Microsoft's
ExCyTIn-Bench~\cite{excytin2025} uses graph-grounded, step-level rewards
tied to a known SOC investigation graph. ExploitBench~\cite{exploitbench2026}
grades along a five-tier, sixteen-capability ladder with a deterministic
verifier and explicitly no LLM-judge. DARPA's AI Cyber
Challenge~\cite{aixcc2025} requires a patch to functionally close a
proven-exploitable vulnerability, not merely to be submitted. None of
these works frame themselves as a response to the outcome-validation
problem in the general case; we argue they are best understood as
convergent, independently arrived-at instances of the same underlying
design principle.

\subsection{The broader CTF and pentesting benchmark landscape}
Beyond the sources scored in Section~\ref{sec:results}, a substantial and
fast-growing body of work evaluates LLM agents on offensive tasks using
the same outcome-only paradigm audited here. HackSynth~\cite{hacksynth2024}
introduces a Planner/Summarizer agent evaluated on 200 PicoCTF and
OverTheWire challenges with binary flag capture. AIRTBench~\cite{airtbench2025}
evaluates frontier models against 70 black-box red-teaming challenges
and reports success rates in the 49 to 61\% range. Fang et
al.~\cite{fang2024oneday} show that GPT-4-class agents can autonomously
exploit real one-day vulnerabilities given a CVE description, using
successful exploitation as the sole outcome signal. Happe and
Cito~\cite{happe2025survey} survey nineteen papers in this space and
find that only around a quarter evaluate open or locally-hosted models,
a coverage gap orthogonal to, but compounding, the grading-design gap
this paper documents. Closer to a method-aware design, PentestJudge~\cite{pentestjudge2025}
proposes hierarchical rubrics that grade an agent's full trajectory
against operational requirements rather than a single terminal check,
and xOffense~\cite{xoffense2025} fine-tunes a Qwen3-32B backbone
specifically for autonomous pentesting and reports 72.72\% on
AutoPenBench. Deng et al.~\cite{deng2025pentest} additionally show that
typed tool interfaces alone shift agent performance by 14 percentage
points independent of the underlying model, a scaffold-level confound
that OVIS treats as orthogonal to, but compounding with, grading-level
integrity.

\subsection{Knowledge and threat-intelligence benchmarks}
A parallel line of work evaluates LLMs on cybersecurity
\emph{knowledge} rather than autonomous exploitation, almost uniformly
using multiple-choice question (MCQ) formats structurally identical to
the outcome-only designs audited here: CTIBench~\cite{ctibench2024} for
threat-intelligence reasoning, SEvenLLM~\cite{sevenllm2024} for
bilingual incident analysis, CyberMetric~\cite{cybermetric2024} for
broad security knowledge sourced from NIST and CWE standards, and Meta's
CyberSecEval~3~\cite{cyberseceval3_2024} for dual-use risk and
capability assessment. The prevalence of MCQ grading across this
adjacent literature suggests the outcome-only pattern we quantify is not
specific to CyberSOCEval or to exploitation benchmarks, but is close to
the default choice across the field, which strengthens the case for a
reusable integrity rubric rather than a one-off critique of any single
benchmark.

\subsection{Meta-synthesis and adjacent security evaluation}
The International AI Safety Report~\cite{aisr2026}, a government and
multi-stakeholder synthesis, itself tracks state-of-the-art trends
across four cyber benchmarks and cites the AIxCC 77\% discovery figure
directly, lending independent weight to the raw capability numbers this
paper treats as inputs rather than as its own contribution. NetSPI's
Open LLM Security Benchmark~\cite{netspi2025} evaluates a related but
distinct property, jailbreak resistance traded off against usability,
using a dual-axis framework that we did not have sufficient
methodological detail to score under OVIS but that shares the same
underlying concern about what a single reported number can and cannot
support.

\subsection{Scaffold and architecture confounds}
A parallel and complementary line of evidence shows that reported
capability is highly sensitive to agent scaffold and tool-schema design
independent of the underlying model~\cite{xbow2026,dcipher2025,enigma2024}.
This paper treats grading-function design as an additional, largely
unexamined confound of the same family, orthogonal to scaffold effects.

\subsection{Production deployment reporting}
Enterprise and hyperscaler deployments of frontier agentic models
against real codebases~\cite{visavvah2026,bigsleep2025,aisle2025} are,
notably, the part of our corpus that most consistently declines to
publish a benchmark-style accuracy number at all, instead reporting
discrete, independently confirmed findings (CVE assignments, maintainer
confirmations) or process metrics such as Visa's Mean Time to
Adapt~\cite{visavvah2026}. We treat this pattern as informative evidence
for RQ2 (Section~\ref{sec:discussion}) rather than a gap in our corpus.

\section{Corpus and Methodology}
\label{sec:method}

\subsection{Corpus construction}
We assembled a corpus of 22 independently published sources spanning
five evidentiary tiers: peer-reviewed or preprint academic benchmarks,
an industry vendor benchmark blog with full per-run transcript
disclosure, a government-run and independently judged competition,
enterprise and hyperscaler production deployments, and an industry
framework self-evaluation. Sources were included only if they reported
at least one quantitative agent performance figure and described, at
some level of detail, how that figure was computed. Of the 22 corpus
sources, \textbf{11} contained sufficient primary-source detail
(published methodology sections, worked examples of the grading
function, or per-run audit narratives) to score all five OVIS dimensions
defensibly; the remaining 11 are retained in the broader corpus for
other analyses in this research program but are excluded from the OVIS
scoring reported here and explicitly flagged as
\textsc{insufficient\_evidence} pending deeper review.

\subsection{The OVIS rubric}
\label{sec:rubric}
Each source $s$ is scored on five dimensions $d \in \{MV, SA, HG, IG,
RT\}$, each on an ordinal scale $\{0,1,2\}$ (absent, partial, explicit
and documented):

\begin{description}[leftmargin=1.4em,itemsep=1pt,topsep=2pt,style=nextline]
\item[\textbf{MV --- Method Verification}] Does the grading function (or
a documented post-hoc audit of it) check \emph{how} the outcome was
reached, rather than only \emph{that} a terminal state was reached?
\item[\textbf{SA --- Shortcut/Contamination Audit}] Is there an explicit,
documented audit for shortcut-solving, answer-leakage, or reward
hacking, with results reported rather than merely asserted?
\item[\textbf{HG --- Human Ground-Truth Anchor}] Is agent performance
anchored to an independent human baseline, such as solve time,
maintainer confirmation, or expert review of the evaluation instrument
itself?
\item[\textbf{IG --- Grader Independence}] Is the grading function
deterministic or otherwise independent of the system under test, i.e.,
not a same-family LLM-judge or self-graded output?
\item[\textbf{RT --- Reproducibility/Transparency}] Are the harness,
grading logic, or raw per-run transcripts made available for independent
re-audit?
\end{description}

\noindent The Outcome-Validation Integrity Score is
\begin{equation}
\mathrm{OVIS}(s) = \frac{\sum_{d \in \{MV,SA,HG,IG,RT\}} \mathrm{score}_d(s)}{10}\times 100\%.
\label{eq:ovis}
\end{equation}

We additionally classify each scored source's grading-function design
into one of three categories, independent of its OVIS score:
\textsc{outcome-only} (binary flag capture, multiple-choice accuracy, or
pass@$k$ on a terminal state), \textsc{process-aware} (tiered capability
ladders, graph-grounded step rewards, or functional/adversarial
validation of an end effect), and \textsc{outcome-only-but-audited} (an
outcome-only grading function whose integrity was subsequently
established or refuted by an independent manual audit, as in the
TrustedSec case). This classification lets us test whether OVIS varies
systematically with \emph{how} a benchmark was designed to measure
success, independent of the raw score any individual benchmark receives.

\subsection{Scoring procedure}
The author read each source's methodology section, worked examples, and
(where available) released harness code, and assigned scores per
dimension with a written justification quoting the specific
methodological evidence relied upon (single-rater scoring in this
preprint; see Section~\ref{sec:threats}). Algorithm~\ref{alg:ovis}
summarizes the scoring and aggregation procedure as implemented.

\begin{figure}[t]
\centering
\fbox{\begin{minipage}{0.95\linewidth}
\small
\textbf{Algorithm 1:} OVIS corpus scoring\\[2pt]
\textbf{Input:} corpus $C$ of candidate sources\\
\textbf{Output:} ranked table $T$, group statistics $G$
\begin{enumerate}[leftmargin=1.2em,itemsep=0pt,topsep=2pt]
  \item $S \leftarrow \{s \in C : \text{depth}(s) \geq \tau_{\text{evidence}}\}$
  \item \textbf{for} each $s \in S$: score $(MV,SA,HG,IG,RT) \in \{0,1,2\}^5$ with justification; assign $\mathrm{design}(s)$
  \item \textbf{for} each $s \in S$: compute $\mathrm{OVIS}(s)$ per Eq.~\ref{eq:ovis}
  \item $T \leftarrow$ sort $S$ by $\mathrm{OVIS}(s)$ descending
  \item \textbf{for} each design category $k$: $G_k \leftarrow$ mean$(MV)$, mean$(\mathrm{OVIS})$ over $\{s : \mathrm{design}(s)=k\}$
  \item \textbf{return} $T, G$
\end{enumerate}
\end{minipage}}
\caption{OVIS scoring and aggregation procedure.}
\label{alg:ovis}
\end{figure}

\section{Results}
\label{sec:results}

\subsection{Ranked OVIS scores}
Table~\ref{tab:ovis} reports OVIS scores for all 11 scored sources,
ranked descending. Scores range from 30.0\% (CAIBench's Cybench-subset
evaluation, pure pass@1 with deterministic but outcome-only grading) to
90.0\% (tied between TrustedSec's audited local-model benchmark and
DARPA's AIxCC). The corpus mean is $60.0\%$ ($\sigma = 20.0$).

\begin{table*}[t]
\centering
\small
\caption{OVIS scores by source, ranked descending. MV/SA/HG/IG/RT are 0--2 dimension scores; OVIS is the percentage of Eq.~\ref{eq:ovis}.}
\label{tab:ovis}
\begin{tabular}{clcccccccl}
\toprule
Rank & Source & MV & SA & HG & IG & RT & Raw/10 & OVIS (\%) & Design category \\
\midrule
1 & TrustedSec (2026), Juice Shop local-model eval~\cite{trustedsec2026} & 2 & 2 & 1 & 2 & 2 & 9 & 90.0 & outcome-only-but-audited \\
2 & DARPA AIxCC 2025 Finals~\cite{aixcc2025} & 2 & 2 & 1 & 2 & 2 & 9 & 90.0 & process-aware \\
3 & Sultan (2026), XBOW-104 scaffold study~\cite{xbow2026} & 1 & 2 & 0 & 2 & 2 & 7 & 70.0 & outcome-only \\
4 & Visa VVAH / Project Glasswing~\cite{visavvah2026} & 2 & 1 & 2 & 2 & 0 & 7 & 70.0 & process-aware-unpublished \\
5 & Cybench 2024~\cite{cybench2024} & 1 & 0 & 2 & 2 & 2 & 7 & 70.0 & outcome-only \\
6 & ExploitBench 2026 (v8-bench)~\cite{exploitbench2026} & 2 & 1 & 0 & 2 & 1 & 6 & 60.0 & process-aware \\
7 & CyberSOCEval 2025~\cite{cybersoceval2025} & 0 & 1 & 0 & 2 & 2 & 5 & 50.0 & outcome-only \\
8 & Microsoft ExCyTIn-Bench 2025~\cite{excytin2025} & 2 & 0 & 0 & 2 & 1 & 5 & 50.0 & process-aware \\
9 & EnIGMA 2024~\cite{enigma2024} & 0 & 0 & 0 & 2 & 2 & 4 & 40.0 & outcome-only \\
10 & D-CIPHER 2025~\cite{dcipher2025} & 0 & 0 & 0 & 2 & 2 & 4 & 40.0 & outcome-only \\
11 & CAIBench 2025 (Cybench subset)~\cite{caibench2025} & 0 & 0 & 0 & 2 & 1 & 3 & 30.0 & outcome-only \\
\bottomrule
\end{tabular}
\end{table*}

\begin{figure}[t]
\centering
\includegraphics[width=\linewidth]{fig1_ovis_ranked.pdf}
\caption{Ranked OVIS scores by source and grading-design category. Dashed line marks the corpus mean (60.0\%).}
\label{fig:ovis_bar}
\end{figure}

\subsection{Metric design predicts method-verification integrity}
Grouping sources by grading-design category (Table~\ref{tab:groups})
shows a large gap in the method-verification (MV) sub-score between
\textsc{outcome-only} designs ($n{=}6$: CyberSOCEval, Cybench, EnIGMA,
D-CIPHER, CAIBench, and Sultan's XBOW study) and \textsc{process-aware}
designs ($n{=}4$: ExCyTIn-Bench, ExploitBench, AIxCC, Visa VVAH). Mean
MV as a percentage of its 2-point maximum is $16.7\%$ for outcome-only
designs versus $100.0\%$ for process-aware designs; every process-aware
source in our corpus scored the maximum on method verification, while
five of six outcome-only sources scored zero. Mean overall OVIS is
correspondingly higher for process-aware designs ($67.5\%$) than
outcome-only designs ($50.0\%$), though the gap narrows at the OVIS
level because outcome-only designs in this corpus tend to compensate
with strong human ground-truth anchoring (Cybench) or reproducibility
(EnIGMA, D-CIPHER).

\begin{table}[t]
\centering
\small
\caption{Mean MV sub-score and mean OVIS by grading-design category.}
\label{tab:groups}
\begin{tabular}{lccc}
\toprule
Design category & $n$ & Mean MV (\%) & Mean OVIS (\%) \\
\midrule
Outcome-only & 6 & 16.7 & 50.0 \\
Process-aware & 4 & 100.0 & 67.5 \\
Outcome-only, audited & 1 & 100.0 & 90.0 \\
\bottomrule
\end{tabular}
\end{table}

\subsection{Auditing rescues, but is rare}
The single \textsc{outcome-only-but-audited} source in our corpus
(TrustedSec) achieves the joint-highest OVIS score in the entire corpus
(90.0\%) despite using the same grading paradigm (string-match on an
HTTP response body) as five of the six lowest-scoring sources. The
difference is entirely attributable to SA and MV: TrustedSec's authors
decoded every passing JWT and traced every SQLi login across 4{,}800
runs, discovering that a challenge with a nominal 98.5\% pass rate had
an intended-method rate of \textbf{0\%}: across 681 total runs on this
challenge, all six benchmarked models combined, the RS256-to-HS256
algorithm-confusion attack the challenge was designed to test was never
performed. Within the one model whose passing runs were fully decoded
and published (nemotron-3-super, $n{=}69$ passes), every pass was
attributable instead to token replay via an unrelated SQL-injection
exploit (28 of 69) or to the structurally simpler \texttt{alg:none}
forgery attack (41 of 69). No other source in our corpus reports an
equivalently granular per-run forensic audit of its own passing results.

\begin{figure*}[t]
\centering
\includegraphics[width=0.72\linewidth]{fig2_jwt_case.pdf}
\caption{Nominal pass rate versus intended-method rate, TrustedSec \texttt{jwt-weak-medium} case study~\cite{trustedsec2026}. Left bar: full mechanism breakdown for the one model with a fully audited, published per-run decode (nemotron-3-super, $n{=}69$ passing runs); the intended RS256$\to$HS256 exploit was used in 0 of 69. Right bar: the challenge's nominal pass rate as reported across all six benchmarked models.}
\label{fig:jwt_case}
\end{figure*}

\section{Discussion}
\label{sec:discussion}

\subsection{RQ1: Does grading design predict gameability?}
Our results support a qualified yes. Every process-aware source in this
corpus scored the maximum method-verification sub-score; every
outcome-only source scored at or near zero on the same dimension unless
independently audited. This is consistent with, and quantifies, the
mechanism predicted by the reward-hacking literature reviewed in
Section~\ref{sec:method}: a grading function that checks a terminal
string or state cannot, by construction, distinguish an intended
exploitation path from an unintended one that happens to produce the
same terminal artifact, which is precisely the ``proxy versus true
objective'' gap Skalse et al.~\cite{skalse2022reward} formalize and Pan
et al.~\cite{pan2022reward} show gets \emph{worse}, not better, as
agents become more capable. XBOW's flag-shortcut audit and TrustedSec's
JWT-decoding audit are two independent instances of researchers
recognizing this limitation and adding a compensating control outside
the automated grading loop, rather than redesigning the grading function
itself.

\subsection{RQ2: Why do production deployments not report benchmark-style numbers?}
The corpus's two production-deployment sources, Visa VVAH and, by
report, Google's Big Sleep and CodeMender
programs~\cite{bigsleep2025}, are the only entries that explicitly and
voluntarily decline to publish an accuracy percentage of any kind,
instead reporting either discrete, externally confirmed findings (CVE
assignments, maintainer acknowledgment) or a purpose-built process
metric. Visa's Mean Time to Adapt explicitly measures inventory
freshness, exploitable paths per release, and validation-cycle time
rather than a closure-count or detection-rate. We interpret this as
indirect but convergent evidence for our central claim: organizations
whose findings face real adversarial and regulatory scrutiny appear to
have independently concluded that outcome-rate metrics are not
defensible enough to publish, and have engineered around the same
failure mode this paper formalizes, via adversarial validation panels
(Visa's S6/S11 stages) rather than via a rubric. Visa's architecture is,
in OVIS terms, strong on MV, HG, and IG (2, 2, 2) but scores zero on RT,
since none of its findings, transcripts, or grading logic are released
for independent re-audit. This illustrates that process rigor and
public auditability are orthogonal properties: a deployment can be
simultaneously highly rigorous and entirely unverifiable by outside
parties.

\subsection{Practical recommendation}
We recommend that any reported agentic cybersecurity benchmark number be
accompanied by, at minimum, an explicit statement of which OVIS
dimensions were satisfied, particularly MV and SA. Where a grading
function is unavoidably outcome-only, for example because ground-truth
mechanism labeling is expensive to construct at scale, we recommend a
TrustedSec-style sampled forensic audit of a subset of passing runs as a
low-cost, high-value compensating control, and that the audited subset
and its findings be published alongside the headline number.

\section{Threats to Validity}
\label{sec:threats}

\textbf{Corpus coverage and scoring subjectivity.} Eleven of 22
candidate sources were excluded from scoring for insufficient documented
methodological detail. This is a deliberate precision-over-recall choice
but means the scored corpus is not a random or complete sample of the
field. Scoring in this preprint was performed by a single author;
inter-rater reliability across independent scorers has not yet been
established and is planned future work (Section~\ref{sec:future}).

\textbf{Ordinal, not interval, dimensions.} The 0/1/2 scale per
dimension is ordinal; treating OVIS as a linear percentage assumes equal
spacing between ``absent,'' ``partial,'' and ``explicit,'' which may not
hold uniformly across dimensions.

\textbf{Small $n$ per design category.} With $n{=}6$ outcome-only and
$n{=}4$ process-aware sources, the group means in Table~\ref{tab:groups}
are directionally suggestive rather than statistically confirmatory. No
significance test is reported given the sample size, and we explicitly
avoid overclaiming statistical significance here.

\textbf{Self-selection in what gets published in this level of detail.}
Sources with the methodological detail needed to score OVIS at all may
be systematically more rigorous authors to begin with, which could
inflate the corpus's overall integrity relative to the true population
of agentic cybersecurity benchmarks in circulation, including many
vendor leaderboards with no published methodology at all.

\section{Future Work}
\label{sec:future}
Immediate extensions include expanding the scored corpus toward the full
22-source pool and beyond as methodological detail is obtained for the
currently-excluded sources, multi-rater scoring with reported
inter-rater agreement such as Cohen's $\kappa$, and applying the OVIS
rubric outside cybersecurity to agentic coding and computer-use
benchmarks that share the same terminal-state grading pattern. A
further direction, informed by the psychometrics literature on item
response theory (IRT) for NLP evaluation~\cite{lalor2016irt}, is to fit
a joint model of benchmark item difficulty and model ability across
sources, conditioned on OVIS as a covariate, to test whether low-OVIS
benchmarks systematically overstate capability relative to high-OVIS
benchmarks measuring comparable tasks.

\section{Conclusion}
We introduced the Outcome-Validation Integrity Score, a five-dimension
rubric for auditing whether an agentic cybersecurity benchmark's grading
function can be satisfied without the intended capability being
exercised. Applied to eleven independently documented sources, OVIS
shows a substantial gap in method-verification integrity between
outcome-only grading designs (16.7\%) and process-aware designs
(100.0\%), and identifies post-hoc manual auditing as a rare but
effective compensating control, illustrated by a documented case in
which a nominally 98.5\%-passing benchmark challenge had a true
intended-method success rate of zero across 681 runs. We release our
rubric and corpus to support future, larger-scale, multi-rater audits of
the rapidly growing body of agentic cybersecurity evaluation claims.

\begin{thebibliography}{99}

\bibitem{xbow2026} Z. Sultan, ``Backbone or Backbone-Architecture? A
Controlled Study of LLM Agents on Web Penetration-Testing CTFs,''
working draft, Jun. 2026.

\bibitem{cybersoceval2025} L. Deason \textit{et al.}, ``CyberSOCEval:
Benchmarking LLMs Capabilities for Malware Analysis and Threat
Intelligence Reasoning,'' Meta / CrowdStrike, Sep. 2025, rev. Nov. 2025.

\bibitem{excytin2025} Microsoft Security, ``ExCyTIn-Bench: Evaluating
LLM Agents on Autonomous Threat Investigation,'' Microsoft Security
Blog, 2025.

\bibitem{exploitbench2026} ExploitBench Contributors, ``ExploitBench:
Tiered Capability Evaluation for LLM Exploit Synthesis (v8-bench),''
2026.

\bibitem{trustedsec2026} B. McGrath, ``Benchmarking Self-Hosted LLMs for
Offensive Security,'' \textit{TrustedSec Blog}, Apr. 2026.

\bibitem{visavvah2026} R. Taneja and S. Kumaraswamy, ``Visa Releases Its
AI-Powered Cyber Defense System to Open Source,'' Visa Corporate
Perspectives, Jun. 2026; Visa, ``Visa Vulnerability Agentic Harness,''
GitHub, 2026.

\bibitem{aixcc2025} DARPA / ARPA-H, ``AI Cyber Challenge (AIxCC) Final
Competition Results,'' Aug. 2025.

\bibitem{cybench2024} A. K. Zhang \textit{et al.}, ``Cybench: A
Framework for Evaluating Cybersecurity Capabilities and Risks of
Language Models,'' \textit{arXiv:2408.08926}, 2024.

\bibitem{enigma2024} T. Abramovich \textit{et al.}, ``EnIGMA:
Interactive Tools Substantially Assist LM Agents in Finding Security
Vulnerabilities,'' \textit{arXiv:2409.16165}, 2024.

\bibitem{dcipher2025} ``D-CIPHER: Dynamic Collaborative Intelligent
Multi-Agent System with Planner and Heterogeneous Executors for
Offensive Security,'' \textit{arXiv:2502.10931}, 2025.

\bibitem{caibench2025} Alias Robotics, ``Cybersecurity AI Benchmark
(CAIBench): A Meta-Benchmark for Evaluating Cybersecurity AI Agents,''
\textit{arXiv:2510.24317}, 2025.

\bibitem{deng2025pentest} G. Deng \textit{et al.}, ``What Makes a Good
LLM Agent for Real-world Penetration Testing?,'' \textit{arXiv:2602.17622},
2026.

\bibitem{cai2025} Alias Robotics, ``CAI: Cybersecurity AI, an Open
Framework for Autonomous Offensive and Defensive Security,'' GitHub,
2025.

\bibitem{bigsleep2025} Google DeepMind and Google Project Zero, ``Big
Sleep: Real-World Vulnerability Discovery with Large Language Models,''
Google Security Blog, Nov. 2024; ``Google AI `Big Sleep' Stops
Exploitation of Critical SQLite Vulnerability Before Hackers Act,''
\textit{The Hacker News}, Jul. 2025.

\bibitem{aisle2025} AISLE, ``AISLE Discovered New OpenSSL
Vulnerabilities,'' AISLE Research Blog, Oct. 2025.

\bibitem{fuzzingbrain2026} Z. Sheng, Z. Chen, Q. Xu, K. Zhu, and J.
Huang, ``FuzzingBrain V2: A Multi-Agent LLM System for Automated
Vulnerability Discovery and Reproduction,'' \textit{arXiv:2605.21779},
2026.

\bibitem{amodei2016concrete} D. Amodei, C. Olah, J. Steinhardt, P.
Christiano, J. Schulman, and D. Man\'e, ``Concrete Problems in AI
Safety,'' \textit{arXiv:1606.06565}, 2016.

\bibitem{skalse2022reward} J. Skalse, N. H. R. Howe, D. Krasheninnikov,
and D. Krueger, ``Defining and Characterizing Reward Hacking,''
\textit{arXiv:2209.13085}, 2022.

\bibitem{pan2022reward} A. Pan, K. Bhatia, and J. Steinhardt, ``The
Effects of Reward Misspecification: Mapping and Mitigating Misaligned
Models,'' \textit{arXiv:2201.03544}, 2022.

\bibitem{krakovna2020specification} V. Krakovna \textit{et al.},
``Specification Gaming: The Flip Side of AI Ingenuity,'' DeepMind
Research Blog, Apr. 2020.

\bibitem{hacksynth2024} L. Muzsai, D. Imolai, and A. Luk\'acs,
``HackSynth: LLM Agent and Evaluation Framework for Autonomous
Penetration Testing,'' \textit{arXiv:2412.01778}, 2024.

\bibitem{happe2025survey} A. Happe and J. Cito, ``Benchmarking Practices
in LLM-Driven Offensive Security: Testbeds, Metrics, and Experiment
Design,'' \textit{arXiv:2504.10112}, 2025.

\bibitem{airtbench2025} ``AIRTBench: Measuring Autonomous AI Red
Teaming Capabilities in Language Models,'' \textit{arXiv:2506.14682}, 2025.

\bibitem{pentestjudge2025} ``PentestJudge: Judging Agent Behavior
Against Operational Requirements,'' \textit{arXiv:2508.02921}, 2025.

\bibitem{xoffense2025} ``xOffense: AI-driven Autonomous Pentesting with
Offensive Knowledge-Enhanced LLMs,'' \textit{arXiv:2509.13021}, 2025.

\bibitem{fang2024oneday} R. Fang, R. Bindu, A. Gupta, and D. Kang,
``LLM Agents can Autonomously Exploit One-day Vulnerabilities,''
\textit{arXiv:2404.11625}, 2024.

\bibitem{ctibench2024} M. T. Alam, D. Bhusal, L. Nguyen, and N. Rastogi,
``CTIBench: A Benchmark for Evaluating LLMs in Cyber Threat
Intelligence,'' \textit{arXiv:2406.07599}, 2024.

\bibitem{sevenllm2024} H. Ji \textit{et al.}, ``SEvenLLM: Benchmarking,
Eliciting, and Enhancing Abilities of Large Language Models in Cyber
Threat Intelligence,'' \textit{arXiv:2405.03446}, 2024.

\bibitem{cybermetric2024} N. Tihanyi, M. A. Ferrag, R. Jain, T. Bisztray,
and M. Debbah, ``CyberMetric: A Benchmark Dataset Based on
Retrieval-Augmented Generation for Evaluating LLMs in Cybersecurity
Knowledge,'' \textit{arXiv:2402.07688}, 2024.

\bibitem{cyberseceval3_2024} S. Wan \textit{et al.}, ``CyberSecEval 3:
Advancing the Evaluation of Cybersecurity Risks and Capabilities in
Large Language Models,'' \textit{arXiv:2408.01605}, 2024.

\bibitem{aisr2026} ``International AI Safety Report 2026,''
\textit{arXiv:2602.21012}, 2026.

\bibitem{netspi2025} NetSPI, ``Open LLM Security Benchmark: Balancing
Security and Usability of Large Language Models,'' NetSPI, 2025.

\bibitem{lalor2016irt} J. P. Lalor, H. Wu, and H. Yu, ``Building an
Evaluation Scale using Item Response Theory,'' in \textit{Proc. EMNLP},
2016, \textit{arXiv:1605.08889}.

\end{thebibliography}

\end{document}
